\begin{ExerciseBlock}
\ExerciseHeading{}{Verständnisfragen}{forces_checks}
Entscheiden Sie, welche Aussagen wahr oder falsch sind, und begründen Sie Ihre Entscheidung.
\begin{enumerate}
\item Es gibt 2 Newtonsche Axiome der Mechanik.\TrueFalseAnswer{falsch}
\item Die Kraft entspricht der zeitlichen Ableitung des Impulses.\TrueFalseAnswer{wahr}
\item Ein Zug fährt auf der Erde von Süden nach Norden. Die Corioliskraft, die auf einen fahrenden Zug einwirkt, ist am Äquator am stärksten.\TrueFalseAnswer{falsch}
\item Die Zentrifugalkraft infolge der Erdrotation ist für gleiche Massen auf der Erdoberfläche am Äquator am größten.\TrueFalseAnswer{wahr}
\item Bei einer idealen Feder verdoppelt sich der Betrag der Rückstellkraft, wenn sich die Auslenkung verdoppelt.\TrueFalseAnswer{wahr}
\item Für eine vorgegebene Richtung lässt sich $\vec F$ eindeutig in eine parallele und eine senkrechte Komponente zerlegen.\TrueFalseAnswer{wahr}
\item Eine große Kugel hat einen größeren $c_W$-Wert als eine kleine Kugel.\TrueFalseAnswer{falsch}
\item Aktionsprinzip: Für eine konstante Kraft halbiert sich die Beschleunigung bei Verdopplung der Masse.\TrueFalseAnswer{wahr}
\item Das Hook'sche Gesetz gilt nur für unbelastete Federn.\TrueFalseAnswer{falsch}
\item Die Haftreibungskraft ist immer größer als die Gleitreibungskraft.\TrueFalseAnswer{falsch}
\item Das Stokes'sche Reibungsgesetz gilt, wenn sich hinter dem Objekt Wirbel ausbilden.\TrueFalseAnswer{falsch}
\item Scheinkräfte treten in Inertialsystemen niemals auf.\TrueFalseAnswer{wahr}
\item Corioliskräfte treten in geradlinig beschleunigten Bezugssystemen auf.\TrueFalseAnswer{falsch}
\item Zu jeder Scheinkraft gibt es eine Gegenkraft.\TrueFalseAnswer{falsch}
\item Je härter eine Feder, desto größer ist die Federkonstante $D$.\TrueFalseAnswer{wahr}
\item Jede Reibungskraft wirkt stets antiparallel zur Geschwindigkeit eines Körpers.\TrueFalseAnswer{falsch}
\item Die inneren Kräfte beeinflussen den Gesamtimpuls eines Teilchensystems.\TrueFalseAnswer{falsch}
\end{enumerate}
\end{ExerciseBlock}
